% Generated draft: requires mathematical and source review.
\documentclass[11pt]{article}
\usepackage[margin=25mm]{geometry}
\usepackage{fontspec}
\setmainfont{Latin Modern Roman}
\usepackage{amsmath,amssymb,mathtools,unicode-math}
\setmathfont{Latin Modern Math}
\usepackage{enumitem,needspace,xurl,hyperref}
\setlength{\parindent}{0pt}
\setlength{\parskip}{5pt}
\setlength{\emergencystretch}{3em}
\newcommand{\N}{\mathbb{N}}
\newcommand{\Z}{\mathbb{Z}}
\newcommand{\Q}{\mathbb{Q}}
\newcommand{\R}{\mathbb{R}}
\newcommand{\C}{\mathbb{C}}
\newcommand{\F}{\mathbb{F}}
\newcommand{\A}{\mathbb{A}}
\newcommand{\PP}{\mathbb{P}}
\newcommand{\E}{\mathbb{E}}
\newcommand{\eps}{\varepsilon}
\newcommand{\dd}{\,\mathrm{d}}
\begin{document}
\section*{UnsolvedMath ID: 20002043}
\textbf{Source problem code:} catalogue entry
\subsection{Definitions and mathematical statement}
Let \((M^n,[g])\) be an oriented Riemannian conformal manifold of even dimension. A Weyl structure is a torsion-free connection \(D\) preserving the conformal class. In a representative metric write \(Dg=-2\alpha\otimes g\); under \(\widehat g=e^{2\omega}g\), \(\widehat\alpha=\alpha-\dd\omega\). Its Faraday form is \(F_D=\dd\alpha\), independent of the representative. Call \(D\) closed if \(F_D=0\), and exact if \(D\) is the Levi-Civita connection of a global representative.

Problem 22, whose antecedent is in companion 20002042, asks for a natural top form \(\mathcal Q([g],D)\), equivalently a scalar density of weight \(-n\), independent of representative and reducing to \(Q_g\,\dd V_g\) when \(D=\nabla^g\). The Riemannian \(Q\)-curvature normalization is part of the choice. The source suggests that existence might restrict that choice. The supplied aid attributes a general natural cotractor/paired-density construction for smooth Weyl structures in all even \(n\geq4\) to Hirachi, Luebbe and Matsumoto, with dimension two handled elementarily; full references and construction conventions are not supplied.

For a closed Weyl structure, the aid gives a patching statement. Locally \(\alpha=\dd\varphi_i\), so \(g_i=e^{2\varphi_i}g\) is \(D\)-parallel. On connected overlaps the \(g_i\) differ by constant homotheties. Since the chosen critical metric \(Q\)-form is unchanged by constant homothety, the forms \(Q_{g_i}\,\dd V_{g_i}\) patch canonically for that choice of metric \(Q\), including nonexact closed structures. This does not require a single globally \(D\)-parallel metric.

The displayed four-dimensional answer/context attributed to Eastwood--Slov\'ak is
\[
\mathcal Q_4=\bigl(2\mathbf P^2-2\mathbf P^{ab}\mathbf P_{ba}-D_aD^a\mathbf P\bigr)\,\dd V_g.
\]
Here \(\mathbf P_{ab}\) is the Weyl Rho tensor, potentially nonsymmetric, \(\mathbf P=g^{ab}\mathbf P_{ab}\), and the Weyl derivatives act on the appropriate conformal densities. The extract suppresses those weights and much tensor typography; it is not an unweighted scalar formula independent of volume.

The aid reports the convention-sensitive identity
\[
\mathbf P^{ab}\mathbf P_{ba}
=\mathbf P^{ab}\mathbf P_{ab}-\tfrac12F_D^{ab}(F_D)_{ab}.
\]
The reversed contraction is therefore material. Away from the closed locus, \(\mathcal Q+cF_D^{\wedge n/2}\) is another natural top-form extension with the same reduction on metric structures. In dimension four, \(c|F_D|_g^2\,\dd V_g\) is also conformally invariant and can change total integrals on closed manifolds when \(F_D\ne0\). Wedge-power ambiguities alone need not change those totals, since a global representative makes \(F_D\) exact. The aid says the Eastwood--Slov\'ak and Hirachi--Luebbe--Matsumoto extensions can have different total integrals; no full comparison formula is supplied. Existence therefore does not imply uniqueness off the closed locus, even though all such Faraday additions vanish when \(D\) is locally metric.
\subsection{Short English statement}
Extend metric \(Q\)-curvature naturally to Weyl connections in every even dimension. The supplied aid reports a general construction and canonical local-metric patching for closed Weyl structures. Faraday terms show that agreement with metric \(Q\)-curvature does not ensure uniqueness for nonclosed structures.
\subsection{Sources}
\begin{itemize}
\item Problem 22, AIM workshop \emph{Conformal structure in geometry, analysis, and physics}: \url{https://aimath.org/WWN/confstruct/confstruct.pdf}. The following four-dimensional formula is attributed to Eastwood and Slov\'ak.
\item The supplied aid attributes the general extension to Hirachi, Luebbe and Matsumoto, without title, year or identifier; the original literature field is empty.
\item Local registry \texttt{data/unsolvedmath/problems.json}, permanent ID 20002043, code \texttt{catalogue entry}; canonical \texttt{aim-geometry-notes.json}, \texttt{notes[380]}; updated 14 August 2026. Antecedent companion ID 20002042, code \texttt{AIM-GEOMETRY-0380}, \texttt{notes[379]}, provides the preceding Weyl-structure extension discussion.
\item Attempt 1: partial result, problem status at run solved, rigor label proved, novelty confidence low; current solved and solved-in-literature labels. The machine-generated aid is not peer reviewed and requires independent verification before citation.
\end{itemize}
\end{document}
